% Expository recovery of INTERRELACIONES_ESTRUCTURALES_VACIO.md, §§6--8.
% Insert after cubic inversion and before logarithmic coordinates.
\subsection{Transported composition of constitutive responses}
\label{amp-comp-sec}

The angular transport and the incidences $90$ and $120$ originate in the
same APP--TRIT--TPK structure: the sheets retain evaluations and their
carries; the tritic regime retains orientation; the TPK calendar transports
phase and memory to the angular channels. The constitutive response
applies the function $f$ of \eqref{iii:eq:monotone} to those channels.
Its inversion expresses, in response coordinates, the composition of
transports sharing a common frame. The constructive order is thereby
preserved: transport, composition and spectral evaluation. The resulting
operation acts on this representation of the enriched state; its inverse
recovers the channels, while the complete memory history remains in the
underlying structure.

Let $s=q^{30}$ and $\psi=f^{-1}$. The exponent $30$ is the greatest
common divisor of the incidences $90$ and $120$. Cubic inversion determines
$\psi:(3/4,1)\longrightarrow(0,1)$. Adjoining the boundary response
$f(1)=3/4$ extends the bijection to
\[
 \psi:D\longrightarrow(0,1],\qquad D=[3/4,1).
\]
At this boundary the rational function $f$ is used; the quotient
$(1-q^{90})/(1-q^{120})$ has a removable singularity there.

\begin{proposition}[Response monoid and additive coordinate]
\label{amp-comp-monoid}
The operation
\begin{equation}
 r\odot t=f\bigl(\psi(r)\psi(t)\bigr),\qquad r,t\in D,
 \label{amp-comp-law}
\end{equation}
makes $D$ a commutative cancellative monoid with identity $3/4$.
The map
\begin{equation}
 \ell(r)=-\frac1{30}\log\psi(r)
 \label{amp-comp-character}
\end{equation}
is an isomorphism from this monoid to $(\mathbb R_{\geq0},+)$.
The interval $(3/4,1)$ is closed under $\odot$; the identity belongs
to the boundary extension. The identity is the only invertible element
within $D$.
\end{proposition}
\begin{proof}
The product of two elements of $(0,1]$ belongs to the same interval, and
\[
 \psi(r\odot t)=\psi(r)\psi(t).
\]
For three responses, either association has image
$\psi(r)\psi(t)\psi(u)$; injectivity of $\psi$ proves associativity.
Commutativity follows from scalar multiplication. If
$r\odot t=r\odot u$, positivity of $\psi(r)$ permits cancellation,
giving $\psi(t)=\psi(u)$ and hence $t=u$. The parameter $1$ corresponds
to $3/4$ and provides the identity. The logarithm transforms products into
sums, and its range on $(0,1]$ shows that $\ell$ is a bijection onto
$\mathbb R_{\geq0}$. In particular, the response of a channel $q$
satisfies $\ell(r)=-\log q$. If both parameters are strictly less than
$1$, so is their product. Finally, two numbers in $(0,1]$ have product
$1$ only when both are $1$.
\end{proof}

\begin{proposition}[Admissible cancellation and inverse extension]
\label{amp-comp-inverse}
For $r,u\in D$, the equation $r\odot t=u$ has a solution $t\in D$
exactly when $u\geq r$. This solution is unique and is given by
\begin{equation}
 t=f\!\left(\frac{\psi(u)}{\psi(r)}\right).
 \label{amp-comp-cancellation}
\end{equation}
The same function $f$ is a bijection from $(0,\infty)$ onto $(0,1)$.
With this extension of $\psi$, the law~\eqref{amp-comp-law} transports
the multiplicative group of positive parameters to the interval $(0,1)$.
Group inversion satisfies
\begin{equation}
 r^{\ominus}=\psi(r)r,\qquad
 r\odot r^{\ominus}=\frac34.
 \label{amp-comp-group-inverse}
\end{equation}
\end{proposition}
\begin{proof}
The equation is equivalent to
$\psi(t)=\psi(u)/\psi(r)$. This quotient belongs to $(0,1]$ exactly
when $\psi(u)\leq\psi(r)$, which, by decreasing monotonicity, is
equivalent to $u\geq r$. The bijection proves existence and uniqueness.
The derivative in \eqref{iii:eq:monotone} is negative for every $s>0$;
the limits of $f$ at $0$ and infinity are $1$ and $0$, respectively.
This establishes the extended bijection. Multiplying the numerator and
denominator by $s^3$ gives
\[
 f(s^{-1})=\frac{s^3+s^2+s}{s^3+s^2+s+1}=s f(s).
\]
The multiplicative inverse of $s=\psi(r)$ therefore has response
$\psi(r)r$. Its composition with $r$ corresponds to the parameter $1$.
\end{proof}

The extension to $s>1$, equivalently to $q>1$, is the algebraic extension
of the reader. The chapter's physical contractive sector retains its
declared domain. The coordinate $\ell$ extends to an isomorphism with
$(\mathbb R,+)$ and changes sign under group inversion.

\subsubsection{Operator realization in a common sheet frame}

\begin{proposition}[Composition with marked projectors]
\label{amp-comp-operator}
Let $\mathcal R=\mathcal R^*$, $\mathcal R^2=I$ and
$P_\pm=(I\pm\mathcal R)/2$. For $j=a,b$, let
\[
 T_j=q_{j,+}P_++q_{j,-}P_-,\qquad
 \mathcal V_j=f(T_j^{30}),\qquad 0<q_{j,\pm}<1.
\]
The composition $T_{a\circ b}=T_aT_b$ satisfies
\begin{equation}
 \mathcal V_{a\circ b}
 =f\bigl(\psi(\mathcal V_a)\psi(\mathcal V_b)\bigr).
 \label{amp-comp-operator-law}
\end{equation}
Its two responses are $r_{a,\pm}\odot r_{b,\pm}$. If
$T_j=\exp[-(x_jI+y_j\mathcal R)]$ with $x_j>|y_j|$, then
\begin{equation}
 T_aT_b=
 \exp[-((x_a+x_b)I+(y_a+y_b)\mathcal R)],
 \label{amp-comp-angular-addition}
\end{equation}
and $x_a+x_b>|y_a+y_b|$.
\end{proposition}
\begin{proof}
Orthogonality $P_+P_-=0$ gives
\[
 T_aT_b=q_{a,+}q_{b,+}P_++q_{a,-}q_{b,-}P_-.
\]
In particular, $T_a$ and $T_b$ commute. Functional calculus on these
projectors yields
$\psi(\mathcal V_j)=T_j^{30}$ and
$(T_aT_b)^{30}=T_a^{30}T_b^{30}$; applying $f$ proves
\eqref{amp-comp-operator-law}. More explicitly, for every function $g$
defined on the two eigenvalues,
\[
 P_\pm g(T_j)=g(q_{j,\pm})P_\pm.
\]
Since $q_{j,\pm}=\exp[-(x_j\pm y_j)]$, their products prove
\eqref{amp-comp-angular-addition}. The triangle inequality concludes
$x_a+x_b>|y_a|+|y_b|\geq|y_a+y_b|$.
\end{proof}

The recovered coordinates can be written directly as
\[
 x=\frac{\ell(r_+)+\ell(r_-)}2,\qquad
 y=\frac{\ell(r_+)-\ell(r_-)}2.
\]
Composition adds these pairs. Simultaneous sheet exchange permutes the
two channels of each transport, acts as $(x,y)\mapsto(x,-y)$ and is an
automorphism of the componentwise law. Inversion of both transports in
the algebraic extension acts as $(x,y)\mapsto(-x,-y)$. The two
involutions commute and retain distinct functions: one exchanges
orientation markings; the other inverts the transport.

To represent exchange by conjugation, use an operator $L$ satisfying
$L\mathcal R L^{-1}=-\mathcal R$. In the representation
\[
 \mathcal R=\begin{pmatrix}0&1\\1&0\end{pmatrix},\qquad
 L=\begin{pmatrix}1&0\\0&-1\end{pmatrix},
\]
direct multiplication gives $LP_\pm L^{-1}=P_\mp$, and hence
$LT_jL^{-1}=q_{j,-}P_++q_{j,+}P_-$. Conjugation by $\mathcal R$
preserves its own projectors. A trace combining the channels likewise
cannot replace the projectors in this functional identity: their labels
contribute to the composition.

The common-frame hypothesis is substantive. For example, the positive
definite operators
\[
 A_0=\begin{pmatrix}1/2&0\\0&1/3\end{pmatrix},\qquad
 B_0=\begin{pmatrix}1/2&1/6\\1/6&1/2\end{pmatrix}
\]
have positive eigenvalues, but
\[
 A_0B_0=\begin{pmatrix}1/4&1/12\\1/18&1/6\end{pmatrix}
\]
is not self-adjoint. The proposition specifies transport composition on
common projectors; its domain does not encompass an arbitrary mixture of
material media. Realizing a concatenation of HMT routes additionally
retains their memory data: algebraic closure of the parameter cone
describes the representation and does not by itself select new primary
routes.

\subsubsection{Channel product and transport composition}

The normalized speed of one state remains determined by
\[
 \widehat c=(r_+r_-)^{-1}.
\]
Here the ordinary product combines the two responses of that same state.
The law $\odot$ evaluates a different operation: composition of two
transports on each sheet. The following exact calculation distinguishes
their arguments:
\begin{equation}
 f(1/2)=\frac{14}{15},\qquad
 \frac{14}{15}\odot\frac{14}{15}=f(1/4)=\frac{84}{85}
 \ne\frac{196}{225}=\left(\frac{14}{15}\right)^2.
 \label{amp-comp-rational-example}
\end{equation}
The fractions in this example are algebraic test parameters. The
expression for $\widehat c$ and the preceding constitutive identities
remain unchanged. The added content is an explicit law transporting
composition and recovering its additive angular coordinate.
